\thispagestyle{empty}
\begin{center}
\includegraphics[width=\textwidth]{logos/RJBlogoLongo}

\vspace*{1.25in}
\includegraphics[width=\textwidth]{\carphotoid}

\vspace*{0.4in}
\huge{\textbf{\carmake \ \carmodel \ \caryear}}

\vspace*{0.2in}
\begin{large}
\carintervention

\carowner

\carstartdate \ -\ \carenddate

\end{large}
\end{center}

\newpage

\tableofcontents

\newpage

\listoffigures

\newpage

\section{Introduction}

%This section aims at documenting, by providing pieces of evidence, the intervention undertaken at Raimundo Branco on the identified historic vehicle. 
%este documento tem como objetivo a apresetacao de um conjunto de evidencias

This electronic document is a signed pdf file that can be viewed and printed, but that cannot be tampered with to guarantee the veracity of its contents. The electronic version has clickable links to allow quick access to contents that cannot be printed, namely short videos and diagrams.

According to the \emph{Charter of Turin} (please see \hyperref[sec:charter_of_turin]{Annex A}), \textbf{restoration} is the process of ``\textit{replacing missing parts or areas with the aim of displaying an earlier state of the vehicle and goes further than conservation}''. At Raimundo Branco, we guarantee that restored areas discreetly blend in with the existing historical stock.

A team of researchers from the \href{https://www.unl.pt/}{NOVA University} and \href{https://www.iscte-iul.pt/}{University Institute of Lisbon}, in collaboration with Raimundo Branco, developed a model for classic car body interventions matching the guidelines proposed in the \emph{Charter of Turin}. This model was built using \href{https://en.wikipedia.org/wiki/Business_Process_Model_and_Notation}{\emph{BPMN (Business Process Modelling and Notation)}}, the mainstream graphical notation for depicting the steps in a business process. You can browse interactively through the several diagrams of this model \href{http://194.210.120.34:5000/diagram}{\textbf{\underline{here}}}.

The following subsections detail each task performed on the classic vehicle during the intervention, for which relevant evidence was collected. The whole intervention included many other tasks as can be understood by clicking the link to directly access the BPMN diagram where the concrete activity appears. Therefore, the latter provides a context to understand the scope of the task being described.

For each represented task, besides the corresponding \emph{Charter of Turin} guidelines, whenever available, are reproduced the photos that were taken during it are shown, along with explanatory comments, whenever required. In some cases, the collected evidence also includes short videos, for which access links are provided. For getting permission to view the videos on the \href{https://www.pinterest.com/}{\emph{Pinterest}} platform, click \href{\carpinterestaccess}{\textbf{\underline{here}}}.

\newpage

\section{Registry Data}

\begin{table}[H]

    \centering
    \Large
    
    \begin{tabular}{|l|p{6cm}|}
        \hline
        CHARACTERISTIC          & VALUE\\
        \hline
        Make                    & \carmake\\ 
        \hline
        Model                   & \carmodel\\ 
        \hline
        Year of Manufacture     & \caryear \\
        \hline
        Licence Plate           & \carlicenceplate\\
        \hline
        Country                 & \carcountry\\ 
        \hline
        Chassis/Frame Number    & \carchassisnumber \\
        \hline
        Engine Make and Number  & \carenginenumber\\
        \hline
    
    \end{tabular}
\end{table}